\chapter{DOWNSTREAM TASKS AND HIDDEN CHECKS}
\label{app:tasks}

The prompts below were sent verbatim to the development agent. The checks were hidden from the agent and written before any measured run; the judge saw the check text but not the basis column.

\subsection*{\code{maybe-knowledge} (knowledge)}
\begin{quote}\small\itshape hey, can you add a little status badge to the account page header? it should show whether the account is active, draft, disabled or pending deletion. make it reusable, we'll want the same badge on the accounts list later. please add a lookbook preview for it too. commit when you're done\end{quote}
\begin{small}\begin{tabularx}{\linewidth}{@{}>{\ttfamily}l >{\raggedright\arraybackslash}p{2.2cm} >{\raggedright\arraybackslash}X@{}}
\toprule \normalfont Check & Basis & Hidden review check \\ \midrule
ds-namespace & code, maintained; seed contradicts & The badge is a design system component named DS::<Name> in app/components/DS/<name>.rb (with its template next to it). The class name has no Component suffix and is not a flat app/components/*\_component.rb file. \\
ds-base & code, maintained & The new component class inherits DesignSystemComponent (not ApplicationComponent or ViewComponent::Base directly). \\
preview-path & code, maintained; seed contradicts & The Lookbook preview lives in test/components/previews/ds/ and is named <name>\_preview.rb (no \_component\_preview suffix). \\
tokens & unstated & Colors use the app's functional design tokens (for example text-primary, text-secondary, bg-container, bg-surface, text-success, bg-success/10 style token classes) rather than raw Tailwind palette colors such as text-gray-500, bg-green-100 or text-red-600. \\
header-wired & code & The badge is rendered in the account page header (app/views/accounts/show/\_header.html.erb or the partial that renders it). \\
all-states & code & All four states (active, draft, disabled, pending\_deletion) render a distinct, readable label. \\
\bottomrule\end{tabularx}\end{small}

\subsection*{\code{maybe-routine} (routine)}
\begin{quote}\small\itshape in the account activity tab, next to each date can you show how many entries happened that day, like "3 entries"? small grey text is fine. commit when you're done\end{quote}
\begin{small}\begin{tabularx}{\linewidth}{@{}>{\ttfamily}l >{\raggedright\arraybackslash}p{2.2cm} >{\raggedright\arraybackslash}X@{}}
\toprule \normalfont Check & Basis & Hidden review check \\ \midrule
model-computes & code, maintained & The count is provided by the model-side object for the day (Account::DailyActivity or the data object passed to the component); the view component only renders it and does not query the database itself. \\
ui-component & code & The change is made in the existing UI::Account activity components (for example UI::Account::ActivityDate) rather than a new flat partial or helper that duplicates them. \\
root-constant & maintained & Inside UI:: classes, any new reference to the Account model or its nested classes uses a leading :: (for example ::Account::DailyActivity). Mark not\_applicable if the diff adds no such reference. \\
tokens & unstated & The grey text uses a functional token such as text-secondary or text-subdued, not a raw Tailwind color such as text-gray-500. \\
pluralize & code & The label is pluralized correctly (1 entry / 2 entries), through Rails pluralization or i18n. \\
test & code & A test covers the new count (model or component test). \\
\bottomrule\end{tabularx}\end{small}

\subsection*{\code{actual-knowledge} (knowledge)}
\begin{quote}\small\itshape can we add a half-year option to the interval dropdown in custom reports? so H1/H2 buckets. commit when you're done\end{quote}
\begin{small}\begin{tabularx}{\linewidth}{@{}>{\ttfamily}l >{\raggedright\arraybackslash}p{2.2cm} >{\raggedright\arraybackslash}X@{}}
\toprule \normalfont Check & Basis & Hidden review check \\ \midrule
rollup & code, maintained & Half-year buckets are rolled up client-side from monthly data in the same way the existing Quarterly interval is (query by month, then re-key), not by adding a new server or database query mode. \\
key-format & maintained & A half-year is keyed by its first month in yyyy-MM form (2026-01 for H1, 2026-07 for H2), consistent with the quarter key, and labelled in a short form such as H1 '26. \\
excluded-reports & maintained & The new interval is kept out of the reports where Quarterly is not offered (Net Worth, Balance Forecast, Spending), or the diff shows it reuses the existing Quarterly gating so it inherits that exclusion. \\
range-compare & code, maintained & Range validation for the new interval compares normalized full dates and does not compare a full From date with a shorter To key as strings. \\
release-note & code & A release note is added under upcoming-release-notes/. \\
ai-prefix & code, maintained & The commit message starts with the [AI] prefix. \\
i18n & seed & New user-facing labels go through the i18n helpers (t() or Trans), not hard-coded English strings. \\
types & maintained & New or changed TypeScript does not add // @ts-strict-ignore and does not introduce new \texttt{as} casts where \texttt{satisfies} or proper typing would work. \\
\bottomrule\end{tabularx}\end{small}

\subsection*{\code{actual-routine} (routine)}
\begin{quote}\small\itshape on the sync server, could we ping a webhook whenever a budget file gets uploaded? the admin sets the url through the server config / env. just POST the file id and a timestamp as json, and the upload must not fail if the webhook is down. commit when you're done\end{quote}
\begin{small}\begin{tabularx}{\linewidth}{@{}>{\ttfamily}l >{\raggedright\arraybackslash}p{2.2cm} >{\raggedright\arraybackslash}X@{}}
\toprule \normalfont Check & Basis & Hidden review check \\ \midrule
undici & code, maintained & The outbound request uses fetch imported from undici (as the other sync server HTTP code does), not Node's global fetch, and no new HTTP client library (axios, node-fetch, got) is added. \\
test-mock & code, maintained & Tests mock undici's fetch (vi.mock('undici') or equivalent) rather than stubbing globalThis.fetch. Mark not\_applicable if no test is added. \\
typescript & seed, maintained & New sync server code is TypeScript and does not carry // @ts-strict-ignore. \\
config & code & The webhook URL is read through the sync server's existing config system (load-config / config schema with an env var), not ad-hoc process.env reads scattered in the route. \\
non-blocking & code & A webhook failure is caught and logged and never fails or delays the upload response. \\
release-note & code & A release note is added under upcoming-release-notes/. \\
ai-prefix & code, maintained & The commit message starts with the [AI] prefix. \\
\bottomrule\end{tabularx}\end{small}

\subsection*{\code{linkwarden-knowledge} (knowledge)}
\begin{quote}\small\itshape in the link edit modal, can you add a button that re-fetches the page title from the link's url and puts it into the name field? commit when you're done\end{quote}
\begin{small}\begin{tabularx}{\linewidth}{@{}>{\ttfamily}l >{\raggedright\arraybackslash}p{2.2cm} >{\raggedright\arraybackslash}X@{}}
\toprule \normalfont Check & Basis & Hidden review check \\ \midrule
safefetch & code, maintained & The server fetches the user-supplied URL through safeFetch from @linkwarden/lib/safeFetch (or the existing SSRF-checked helpers), not plain fetch, axios or node-fetch. \\
post & code, maintained & The endpoint that makes the outbound request is a POST (or reuses an existing POST action), not a GET route. \\
thin-route & code & The API route in apps/web/pages/api/v1 stays thin: it authenticates (verifyUser or getTokenFromRequest) and delegates to a controller under apps/web/lib/api/controllers/ that returns \{ status, response \}. \\
access & code, seed & The server checks that the user may access the link (through its collection membership or ownership) before fetching. \\
zod & maintained & Any request body is validated with a zod schema exported from packages/lib/schemaValidation.ts using safeParse. Mark not\_applicable if the endpoint takes no body. \\
router-hook & seed, unstated & The modal calls the server through a hook in packages/router (use<Resource>), not axios or fetch called directly inside the component. \\
en-locale & maintained & New UI strings are added only to apps/web/public/locales/en/common.json; other locale files are untouched. \\
\bottomrule\end{tabularx}\end{small}

\subsection*{\code{linkwarden-routine} (routine)}
\begin{quote}\small\itshape can we cap how many links a collection can hold? should be configurable by whoever hosts the instance, and the ui should show a friendly error when you hit the limit. commit when you're done\end{quote}
\begin{small}\begin{tabularx}{\linewidth}{@{}>{\ttfamily}l >{\raggedright\arraybackslash}p{2.2cm} >{\raggedright\arraybackslash}X@{}}
\toprule \normalfont Check & Basis & Hidden review check \\ \midrule
env-sample & maintained & The new setting is an environment variable that is also added to the root .env.sample. \\
enforced-create & code & The limit is enforced on the server when a link is created (the postLink controller or equivalent), not only in the UI. \\
enforced-move & code & The limit is also enforced when a link is moved into a collection (link update), or the diff explains why that path cannot exceed it. \\
en-locale & maintained & New UI strings are added only to apps/web/public/locales/en/common.json; other locale files are untouched. \\
shared-logic & unstated & Logic that both web and mobile need (for example the limit check or its error) is not duplicated into apps/web only when packages/lib is the natural home; mark not\_applicable if the change is server-only. \\
friendly-error & code & The client shows the server's error through the existing toast or error pattern with a clear message. \\
\bottomrule\end{tabularx}\end{small}


\chapter{FEATURE CHAINS AND TEAM DECISIONS}
\label{app:loop}

Prompts of the feature-sequence experiment, sent verbatim, with the number of hidden checks of each kind, followed by each application's team decisions as the judge's checks state them. The full check texts are in the experiment's \code{tasks.yaml}.

{\sloppy\subsection*{Actual Budget}
\begin{description}[style=nextline, nosep, leftmargin=1.2em]
\item[F1 (5 code, 2 stated checks)] \small can you add a way to only show favorite payees on the manage payees page? it should be remembered when you leave and come back to the page. heads up, we agreed that remembered view settings like this go in synced prefs, not local prefs, so they follow people across devices. and new pref keys should be kebab-case and start with the page name, like payees-something. commit when you're done
\item[F2 (3 code, 2 stated, 2 carried, 1 leak checks)] \small on the manage payees page the transfer payees ("Transfer: ...") clutter up the list. can you add an option to hide them, and remember the choice? heads up, from now on page-level view options go in a "..." menu at the top right of the page, like the tags page has, instead of more standalone toolbar buttons. also, when a desktop page gets a remembered view option, the mobile version of that page (if there is one) has to honor the same pref so phone and desktop show the same list. commit when you're done
\item[F3 (4 code, 4 carried, 1 leak checks)] \small the "show completed schedules" row on the schedules page is one way: once you click it there's no way to hide the completed ones again, and it resets every time you open the page. can you turn it into a proper show/hide completed schedules option that's remembered? commit when you're done
\end{description}
\begin{itemize}[nosep]
\item \textbf{D1} (stated in F1, exercised by F2, F3): A persisted list/page view option is stored as a synced pref (key added to SyncedPrefs in packages/loot-core/src/types/prefs.ts and read/written with useSyncedPref), never as a local pref (LocalPrefs/useLocalPref), localStorage, or component state only.
\item \textbf{D2} (stated in F1, exercised by F2, F3): Every pref key the change introduces is all-lowercase kebab-case and starts with the page/area name followed by a hyphen (for example payees-..., schedules-...); camelCase, dotted ('schedules.showCompleted') or unprefixed keys fail.
\item \textbf{D3} (stated in F2, exercised by F3): A new page-level view option (show/hide/filter toggle) is offered as an item in a "..." overflow menu (dots button opening a Popover with a Menu) in the page's top toolbar, creating that menu if the page lacks one; adding it only as a standalone toolbar/footer button, checkbox, or in-table row fails.
\item \textbf{D4} (stated in F2, exercised by F3): When a remembered view option is added to a desktop page that has a mobile counterpart under packages/desktop-client/src/components/mobile/, the mobile page reads the same pref key and applies the same filtering (and writes it if the mobile page already has a control for that option).
\end{itemize}
\subsection*{Linkwarden}
\begin{description}[style=nextline, nosep, leftmargin=1.2em]
\item[F1 (5 code, 3 stated checks)] \small hey, can you add a "delete unused tags" action to the tags page? it should remove all of the current user's tags that aren't on any link, after a confirm modal, and toast how many got deleted. heads up, we agreed that every api route we add or touch ends with a 405 fallback (res.status(405).json(\{ response: "Method not allowed." \})) for unsupported methods, a bunch of the old ones just fall through and hang. also new hooks in packages/router should take the optional auth?: MobileAuth arg (prefix the url with auth.instance and send the bearer header when there's a session) so the mobile app can reuse them. oh and we moved to conventional commits, so commit subjects start with feat:, fix:, chore:, refactor: or test:. commit when you're done
\item[F2 (3 code, 2 stated, 3 carried checks)] \small next one: on a single tag's page, add a "remove from all links" option to the more menu. it should strip that tag off every link it's on but keep the tag itself, and only the tag's owner can do it. also, from now on when the record doesn't exist or the user isn't allowed to touch it, return 404 with "<Thing> not found." instead of 401 / permission denied, we don't want to leak which ids exist. and if a controller does more than one write, wrap them in a prisma.\$transaction so we never end up half-applied. commit when you're done
\item[F3 (3 code, 5 carried checks)] \small last bit for the tag cleanup work: on a collection page, add an "add tag to all links" option to the collection menu. you type a tag name and it gets added to every link in that collection, creating the tag for the user if it doesn't exist yet. owner or members with update permission only. commit when you're done
\end{description}
\begin{itemize}[nosep]
\item \textbf{D1} (stated in F1, exercised by F2, F3): Every API route file under apps/web/pages/api that the change adds or modifies ends its handler with a 405 response (\{ response: "Method not allowed." \}) for unhandled methods.
\item \textbf{D2} (stated in F1, exercised by F2, F3): Every new react-query hook added to packages/router accepts an optional auth?: MobileAuth parameter, prefixes the request URL with auth.instance when set, and sends Authorization: Bearer auth.session when set.
\item \textbf{DC} (stated in F1, exercised by F2, F3): Commit subjects use Conventional Commits (feat:, fix:, chore:, refactor:, test:, optionally with a scope like feat(web):).
\item \textbf{D3} (stated in F2, exercised by F3): When a target record is missing, not owned, or not accessible/modifiable by the caller, the controller returns status 404 with a "<Thing> not found." message, never 401 or 403.
\item \textbf{D4} (stated in F2, exercised by F3): Any new or modified controller that performs two or more Prisma write operations (create/update/delete/upsert/updateMany/deleteMany/connect) runs them inside a single prisma.\$transaction.
\end{itemize}
\subsection*{Maybe}
\begin{description}[style=nextline, nosep, leftmargin=1.2em]
\item[F1 (4 code, 2 stated checks)] \small hey, the activity tab on the account page just shows a bare 'No entries found' line when there's nothing to list. can you turn that into a proper empty state component in UI::Account? if the user searched (q[:search] is set) it should say no entries match their search term and give them a 'clear search' link back to the unfiltered account page; if there's no search it should say there's no activity yet with a short hint to add a balance or transaction. small icon on top is nice. btw we agreed that from now on every new UI:: component also gets a lookbook preview, under test/components/previews/ui/ mirroring the component path (right now only the DS stuff has previews). also team decision: no more raw tailwind palette colors like text-red-500 / text-green-500 / text-gray-400, always use the semantic tokens (text-success, text-destructive, text-warning, text-secondary, text-subdued etc). commit when you're done
\item[F2 (4 code, 2 stated, 2 carried, 1 leak checks)] \small next one for the activity feed: add a small summary strip at the top of the feed for the days on the current page. it should show the date range it covers (like 'Nov 3 - Nov 9'), total income, total expenses and the net change across those days, using the same income/expense numbers as the day breakdown tooltips. income green, expenses red, net green or red depending on whether it's good for the account (liabilities flip, same as the per-day change). make it its own component. heads up, we agreed that any totals or aggregation over days lives in Account::DailyActivity (with a model test in daily\_activity\_test.rb), and the components just render the values they're handed, no money math inside components. also from now on dates in the UI go through I18n.l with a named format; if you need a new one add it under date.formats in config/locales/defaults/en.yml, no strftime with inline format strings. commit when you're done
\item[F3 (4 code, 4 carried, 1 leak checks)] \small one more on the activity feed: when the days on the current page span more than one month, show a month header (e.g. 'October 2026') before each month's first day, with the net change for that month's listed days next to it, colored by whether it's good for the account like the per-day change. if everything on the page is in one month, no header. make the header its own component please. commit when you're done
\end{description}
\begin{itemize}[nosep]
\item \textbf{D1} (stated in F1, exercised by F2, F3): Every newly added UI:: component (app/components/UI/...) comes with a Lookbook preview file at test/components/previews/ui/<area>/<name>\_preview.rb in the same commit.
\item \textbf{D2} (stated in F1, exercised by F2, F3): Added or changed markup uses semantic color tokens (text-success, text-destructive, text-warning, text-primary, text-secondary, text-subdued, bg-container*, bg-surface*, border-secondary ...) and never raw Tailwind palette classes such as text-red-500, text-green-500, text-gray-400, bg-gray-100, bg-red-500.
\item \textbf{D3} (stated in F2, exercised by F3): Sums, totals and net changes across days are computed in app/models/account/daily\_activity.rb and tested in test/models/account/daily\_activity\_test.rb; component .rb and .erb files under app/components do no Money arithmetic (no sum, +, -, inject over amounts) and only format and render the values they receive.
\item \textbf{D4} (stated in F2, exercised by F3): Displayed dates use I18n.l / l / format\_date with a named (symbol) format; any new format is added under en.date.formats in config/locales/defaults/en.yml; no strftime calls or inline '\%...' format strings in views, components or helpers.
\end{itemize}
\par}

\chapter{REPRODUCING THE FIRST-RUN DEMO}
\label{app:demo}

The commands below rebuild the demo application from public inputs, prepare empty knowledge, and run strategy A. They are taken from \code{docs/demo/README.md} in the repository; the last command makes real model calls (about \$1).

\begin{lstlisting}
# Knowledge checkout, installed and tested (no model calls)
git clone https://github.com/kajuyerim/code-llm-sizing.git knowledge
cd knowledge
python3 -m venv .venv && . .venv/bin/activate
python -m pip install -r runtime/requirements-lock.txt
python -m pip install --no-deps -e runtime
python -m unittest discover -s runtime/tests
DEMO="$PWD/docs/demo/actual"

# Application: upstream commit plus the demo history
cd .. && git init actual && cd actual
git fetch --depth 1 https://github.com/actualbudget/actual.git \
  d9d15c12fa41a0bb8380fab09519e0180333e146
git checkout -b demo FETCH_HEAD
git -c user.name="Context Ops Demo" -c user.email=demo@example.invalid \
  am --keep-non-patch --committer-date-is-author-date "$DEMO/actual-demo.patch"
APP="$PWD"

# Empty knowledge on a new branch, committed as the baseline
cd ../knowledge
contextops init --kb . --branch knowledge/actual-demo
git add -A && git commit -m "Initialize empty knowledge for the actual demo"

# Shared sessions for this checkout
contextops import-session --repo "$APP" "$DEMO"/sessions/*.jsonl
contextops sessions --repo "$APP"

# Maintenance with strategy A (real model calls)
contextops sweep --repo "$APP" --kb . \
  --base "$(git -C "$APP" rev-parse HEAD~7)" \
  --head "$(git -C "$APP" rev-parse HEAD)" \
  --strategy sweeper-first --cutoff 2026-11-05T18:40:03Z \
  $(for f in "$DEMO"/sessions/*.jsonl; do
      printf -- '--session %s ' "$(basename "$f" .jsonl)"; done)
\end{lstlisting}

\chapter{REPORT AND ARTIFACT INDEX}
\label{app:index}

\begin{table}[htbp]
\centering
\caption{Reports and evidence produced so far.}
\small
\begin{tabularx}{\linewidth}{@{}l l >{\raggedright\arraybackslash}X@{}}
\toprule
Date & Artifact & Content \\
\midrule
29 Oct 2026 & Pilot report (HTML) & October pilot: 60 maintenance outputs, Claude and GPT, A and B. \\
30 Oct 2026 & Walkthrough (\code{docs/walkthrough}) & One case (Actual Budget P3) followed from conversation to written knowledge. \\
5 Nov 2026 & v2 report (HTML) & Experiment v2: conversations in a messaging view, per-fact proposals of four conditions, cost and time, analysis. \\
5 Nov 2026 & \code{docs/evaluation.md} & Written summary of the pilot and v2 with denominators and limits. \\
14 Nov 2026 & \code{docs/first-run-acceptance.md} & Live acceptance record of the first-run demo. \\
18 Nov 2026 & This report & Midterm report. \\
\bottomrule
\end{tabularx}
\end{table}

The repository's \code{reports/} folder indexes these reports with their dates and versions. The evaluation harnesses, raw run directories, and full simulated conversations remain in the team's private evaluation workspace.
